% Exported from Math++. Compile with pdflatex.
\documentclass[11pt,a4paper]{article}
\usepackage[T1]{fontenc}
\usepackage[utf8]{inputenc}
\usepackage{lmodern}
\usepackage{amsmath,amssymb,amsthm,mathtools}
\usepackage{graphicx,booktabs,xcolor}
\usepackage[hidelinks]{hyperref}
\usepackage[margin=25mm]{geometry}
\providecommand{\degree}{\ensuremath{^\circ}}
\providecommand{\xmark}{\ensuremath{\times}}
\providecommand{\omicron}{o}
\providecommand{\oiint}{\mathop{\oint\mkern-12mu\oint}}
\theoremstyle{plain}
\newtheorem{theorem}{Theorem}
\newtheorem{lemma}[theorem]{Lemma}
\newtheorem{corollary}[theorem]{Corollary}
\newtheorem{proposition}[theorem]{Proposition}
\newtheorem{conjecture}[theorem]{Conjecture}
\theoremstyle{definition}
\newtheorem{definition}[theorem]{Definition}
\newtheorem{example}[theorem]{Example}
\theoremstyle{remark}
\newtheorem*{remark}{Remark}
\setlength{\parindent}{0pt}\setlength{\parskip}{0.6em}
\title{Applications of the integral}
\date{}
\begin{document}
\maketitle
A worksheet of the standard applications. Each operator is written \texttt{\textbackslash{}operatorname\{name\}($\ldots$)}; an interval $[a, b]$ gives the bounds and $x = c$ or $y = c$ gives an axis.

\section*{Area between curves}

The line $y = x$ and the parabola $y = x^2$ cross at $x = 0$ and $x = 1$; without an interval the region they enclose is used.

\begin{equation*}
  \operatorname{Area\ between\ curves}(x, x^2) = \frac{1}{6}
\end{equation*}
Over a whole period of $\sin x$ both humps count as area.

\begin{equation*}
  \operatorname{Area\ between\ curves}(\sin x, [0, 2\pi]) = 4
\end{equation*}
Curves in $y$: $x = y + 2$ and $x = y^2$.

\begin{equation*}
  \operatorname{Area\ between\ curves}(y + 2, y^2) = \frac{9}{2}
\end{equation*}
\section*{Volumes of revolution}

Disks: $y = \sqrt{x}$ on $[0, 4]$ about the $x$-axis.

\begin{equation*}
  \operatorname{Volume\ of\ revolution}(\sqrt{x}, [0, 4]) = 8 \pi
\end{equation*}
Washers about $y = -1$, then shells about the $y$-axis and about $x = 2$, for the region between $y = x$ and $y = x^2$.

\begin{equation*}
  \operatorname{Volume\ of\ revolution}(x, x^2, [0, 1], y = -1) = \frac{7 \pi}{15}
\end{equation*}
\begin{equation*}
  \operatorname{Volume\ of\ revolution}_{shell}(x, x^2, [0, 1], x = 0) = \frac{\pi}{6}
\end{equation*}
\begin{equation*}
  \operatorname{Volume\ of\ revolution}(x, x^2, [0, 1], x = 2) = \frac{\pi}{2}
\end{equation*}
A sphere of radius $2$ from the upper half circle.

\begin{equation*}
  \operatorname{Volume\ of\ revolution}_{disk}(\sqrt{4 - x^2}, [-2, 2]) = \frac{32 \pi}{3}
\end{equation*}
\section*{Lengths and surfaces}

\begin{equation*}
  \operatorname{Arc\ length}(x^{3/2}, [0, 4]) = \frac{80 \sqrt{10}}{27} - \frac{8}{27}
\end{equation*}
\begin{equation*}
  \operatorname{Arc\ length}((\cos t, \sin t), t, [0, 2\pi]) = 2 \pi
\end{equation*}
\begin{equation*}
  \operatorname{Surface\ area}(x, [0, 1]) = \pi \sqrt{2}
\end{equation*}
\section*{Average value, work and centres of mass}

\begin{equation*}
  \operatorname{Average\ value}(x^2, [0, 3]) = 3
\end{equation*}
Stretching a spring with $k = 5$ from $0$ to $0.2$ by Hooke's law $F = kx$.

\begin{equation*}
  \operatorname{Work}(5x, [0, 0.2]) = 0.1
\end{equation*}
\begin{equation*}
  \operatorname{Centroid}(x, x^2, [0, 1]) = \left(\frac{1}{2}, \frac{2}{5}\right)
\end{equation*}
A rod on $[0, 2]$ with density $1 + x$.

\begin{equation*}
  \operatorname{Centre\ of\ mass}(1 + x, [0, 2]) = \frac{7}{6}
\end{equation*}
\section*{The accumulation function}

\begin{equation*}
  \operatorname{Accumulation}(2x, 1) = x^{2} - 1
\end{equation*}
\section*{Improper integrals}

\begin{equation*}
  \int_0^{\infty} \frac{\sin x}{x}\,dx = \frac{\pi}{2}
\end{equation*}
\begin{equation*}
  \int_{-\infty}^{\infty} e^{-x^2}\,dx = \sqrt{\pi}
\end{equation*}
\begin{equation*}
  \int_0^{\infty} \sin(x^2)\,dx \approx 0.6266570687
\end{equation*}
\begin{equation*}
  \int_1^{\infty} \frac{1}{x}\,dx = ?
\end{equation*}
{\small\color{red!60!black}$\times$ $\int$ 1/x dx (near x = $\infty$) diverges}

\end{document}
